\documentclass{notaaula}

        \titulonota{Síntese científica e culminância}
        \creditos{Turma 3B · Ciências da Natureza para o ENEM · Dezembro/2026 · Prof. Josué Cavalcante}

        \begin{document}
        \cabecalho

        \begin{sobre}
        Culminância semestral: organizar evidências, comunicar resultados e propor intervenção responsável em Ciências da Natureza.
        \end{sobre}

        \section{Pergunta, evidência e conclusão}

\begin{copiar}{Estrutura de uma investigação}
Uma síntese científica clara conecta três elementos:
\begin{itens}
\item \dest{Pergunta:} específica, investigável e ligada a um problema real.
\item \dest{Evidência:} dados, medições ou fontes confiáveis, com unidade e incerteza quando cabível.
\item \dest{Conclusão:} resposta limitada pelo alcance dos dados, sem generalização indevida.
\end{itens}
Hipótese é uma explicação testável; não é palpite imune a revisão. Resultado contrário à hipótese
também produz conhecimento quando o método é adequado.
\end{copiar}

\section{Produto de culminância}

\begin{copiar}{Escolha um formato}
O grupo pode produzir infográfico, seminário, podcast científico ou demonstração experimental.
Qualquer formato deve conter:
\begin{itens}
\item problema e relevância para escola ou comunidade;
\item explicação integrada de Física, Química e Biologia;
\item ao menos um dado representado em tabela ou gráfico;
\item fontes identificadas e distinção entre dado, interpretação e opinião;
\item proposta viável, responsável e com forma de avaliar resultado.
\end{itens}
\end{copiar}

\section{Tema-guia: energia na escola}

\begin{exemplo}[Plano de investigação]
\rotulo{Pergunta.} Onde a escola pode reduzir consumo sem prejudicar conforto e aprendizagem?

\rotulo{Dados.} Potência nominal, tempo diário de uso, quantidade de aparelhos e tarifa.
\[
  E_{\mathrm{mês}}=\sum P_i\Delta t_i.
\]
\rotulo{Intervenção.} Comparar iluminação, climatização e hábitos; priorizar ação de maior economia
por custo de implantação. Medir novamente após a mudança.

\rotulo{Limite.} Potência nominal e tempo estimado geram aproximação; clima e ocupação variam.
\end{exemplo}

\begin{copiar}{Do dado à proposta}
Uma proposta de intervenção deve indicar agente, ação, meio, prazo e indicador. Exemplo:
``A equipe de manutenção ajustará horários dos aparelhos durante quatro semanas; o consumo semanal
será comparado à média anterior, corrigindo-se dias sem aula.'' Assim, a ação pode ser avaliada.
\end{copiar}

\section{Comunicação e ética}

\begin{copiar}{Checklist de qualidade}
\begin{itens}
\item Eixos e unidades dos gráficos estão visíveis?
\item Imagens e dados têm fonte?
\item O grupo explicou incertezas e limitações?
\item Pessoas, animais e ambiente foram tratados com segurança e respeito?
\item A conclusão responde à pergunta e não promete além dos dados?
\end{itens}
Plágio, recorte seletivo de resultados e fabricação de dados anulam a confiabilidade do produto.
\end{copiar}

\section{Autoavaliação}

\begin{exercicios}[Roteiro de entrega]
\begin{questoes}
\item Escreva a pergunta central do grupo em uma frase.
\item Liste as três principais grandezas ou variáveis e suas unidades.
\item Identifique variável controlada, independente e dependente, quando aplicável.
\item Produza um gráfico e escreva duas observações que vêm diretamente dos dados.
\item Escreva uma interpretação e deixe claro como ela difere da observação.
\item Registre ao menos duas fontes confiáveis e o motivo da escolha.
\item Descreva uma limitação do método.
\item Formule a intervenção com agente, ação, meio, prazo e indicador.
\item Distribua tarefas e indique responsável por revisão conceitual e revisão visual.
\item Faça um ensaio de três minutos e corte informações que não ajudam a responder à pergunta.
\item Autoavalie de 0 a 2: domínio conceitual, uso de evidências, integração, comunicação e colaboração.
\item Escreva uma aprendizagem que será útil fora da prova do ENEM.
\end{questoes}
\end{exercicios}

\begin{dica}
Um produto curto com pergunta clara, evidência legível e conclusão honesta comunica melhor do que
muitos efeitos visuais sem encadeamento científico.
\end{dica}


\end{document}
